\documentclass[../main.tex]{subfiles}

\begin{document}

\chapter{Optimizer Cheat Sheet}\label{app:cheatsheet}

\begin{center}
\textit{A quick reference for practitioners. Print this and pin it to your wall.}
\end{center}

\vspace{0.5cm}

%% ============================================================
\section{One-Page Optimizer Summary}\label{sec:cheat-optimizer-table}

\begin{table}[htbp]
\centering
\scriptsize
\renewcommand{\arraystretch}{1.3}
\begin{tabularx}{\textwidth}{@{}l>{\raggedright\arraybackslash}p{2.8cm}>{\raggedright\arraybackslash}p{1.8cm}>{\raggedright\arraybackslash}p{2.2cm}>{\raggedright\arraybackslash}p{2.2cm}@{}}
\toprule
\textbf{Optimizer} & \textbf{Update Rule} & \textbf{Defaults} & \textbf{Pros} & \textbf{Cons} \\
\midrule

\textbf{SGD} &
$\theta \leftarrow \theta - \eta g$ &
(required) &
Simple, generalizes well &
Slow, sensitive to LR \\[0.3cm]

\textbf{SGD + Momentum} &
$v \leftarrow \beta v + g$ \newline
$\theta \leftarrow \theta - \eta v$ &
(required) \newline $\beta = 0.9$ &
Faster, smooths noise &
Two hyperparams \\[0.3cm]

\textbf{Nesterov} &
$v \leftarrow \beta v + \nabla f(\theta - \eta\beta v)$ \newline
$\theta \leftarrow \theta - \eta v$ &
(required) \newline $\beta = 0.9$ &
Better convergence &
Slightly complex \\[0.3cm]

\textbf{AdaGrad} &
$G \leftarrow G + g^2$ \newline
$\theta \leftarrow \theta - \frac{\eta}{\sqrt{G + \epsilon}} g$ &
$\eta = 0.01$ \newline $\epsilon = 10^{-8}$ &
Good for sparse data &
LR decays to zero \\[0.3cm]

\textbf{RMSprop} &
$v \leftarrow \alpha v + (1-\alpha) g^2$ \newline
$\theta \leftarrow \theta - \frac{\eta}{\sqrt{v + \epsilon}} g$ &
$\eta = 0.01$ \newline $\alpha = 0.99$ &
Adaptive, non-convex &
Less studied theory \\[0.3cm]

\textbf{Adam} &
$m \leftarrow \beta_1 m + (1-\beta_1) g$ \newline
$v \leftarrow \beta_2 v + (1-\beta_2) g^2$ \newline
$\theta \leftarrow \theta - \frac{\eta \hat{m}}{\sqrt{\hat{v}} + \epsilon}$ &
$\eta = 0.001$ \newline $\beta_1 = 0.9$ \newline $\beta_2 = 0.999$ &
Fast, adaptive, robust &
May generalize worse \\[0.3cm]

\textbf{AdamW} &
Adam + decoupled weight decay \newline
$\theta \leftarrow \theta - \eta\left(\frac{\hat{m}}{\sqrt{\hat{v}} + \epsilon} + \lambda\theta\right)$ &
$\lambda = 0.01$ \newline (others as Adam) &
Best for transformers &
Extra hyperparam \\

\bottomrule
\end{tabularx}
\caption{Optimizer quick reference. Here $g = \nabla_\theta \mathcal{L}$, $\hat{m} = m/(1-\beta_1^t)$, $\hat{v} = v/(1-\beta_2^t)$.}
\label{tab:optimizer-quick-reference}
\end{table}

\begin{tcolorbox}[colback=gray!5, colframe=gray!50, title={\textbf{When to Use Each Optimizer}}, fonttitle=\bfseries]
\begin{itemize}[leftmargin=*, itemsep=2pt]
    \item \textbf{SGD + Momentum}: CNNs for competitions, when you have time to tune, best final accuracy
    \item \textbf{Nesterov}: Same as momentum, slightly faster convergence on convex problems
    \item \textbf{AdaGrad}: Sparse features, NLP with bag-of-words, embeddings with rare tokens
    \item \textbf{RMSprop}: RNNs, non-stationary objectives, online learning
    \item \textbf{Adam}: Quick prototyping, most deep learning tasks, robust default
    \item \textbf{AdamW}: Transformers, BERT/GPT fine-tuning, when using weight decay
\end{itemize}
\end{tcolorbox}

%% ============================================================
\section{Hyperparameter Quick Reference}\label{sec:cheat-hyperparams}

\subsection{Learning Rate by Task}

\begin{table}[htbp]
\centering
\renewcommand{\arraystretch}{1.2}
\begin{tabular}{llll}
\toprule
\textbf{Task} & \textbf{SGD} & \textbf{Adam/AdamW} & \textbf{Notes} \\
\midrule
\textbf{Vision (CNNs)} & 0.1 & 1e-3 to 3e-4 & Decay by 10x at plateaus \\
\textbf{Vision (ViT)} & 0.01 & 1e-4 to 3e-4 & Use warmup \\
\textbf{NLP (Transformers)} & -- & 1e-4 to 5e-5 & Linear warmup essential \\
\textbf{NLP (Fine-tuning)} & -- & 2e-5 to 5e-5 & 10x lower than pretraining \\
\textbf{RL (Policy Gradient)} & 1e-3 & 3e-4 & Sensitive, start low \\
\textbf{RL (Value Function)} & 1e-3 & 1e-3 & Can be higher \\
\textbf{GANs} & -- & 1e-4 to 2e-4 & Often use different LR for G/D \\
\textbf{Quick Prototype} & 0.01 & 3e-4 & The ``just works'' defaults \\
\bottomrule
\end{tabular}
\caption{Recommended learning rate ranges by task and optimizer.}
\label{tab:lr-by-task}
\end{table}

\subsection{Batch Size Guidelines}

\begin{table}[htbp]
\centering
\renewcommand{\arraystretch}{1.2}
\begin{tabular}{lll}
\toprule
\textbf{Scenario} & \textbf{Batch Size} & \textbf{Notes} \\
\midrule
\textbf{Memory-limited GPU} & 16--64 & Use gradient accumulation \\
\textbf{Standard training} & 32--256 & Good bias-variance tradeoff \\
\textbf{Large-scale (with scaling)} & 256--8192 & Scale LR linearly: $\eta' = \eta \cdot B'/B$ \\
\textbf{Transformers} & 256--2048 & Larger often better with warmup \\
\textbf{Fine-tuning} & 16--32 & Smaller preserves pretrained features \\
\textbf{Few-shot / Small data} & 4--16 & Noise can help regularize \\
\bottomrule
\end{tabular}
\caption{Batch size recommendations by scenario.}
\label{tab:batch-size}
\end{table}

\subsection{Other Hyperparameters}

\begin{table}[htbp]
\centering
\renewcommand{\arraystretch}{1.2}
\begin{tabular}{llp{6cm}}
\toprule
\textbf{Hyperparameter} & \textbf{Typical Range} & \textbf{Guidance} \\
\midrule
\textbf{Momentum} ($\beta$) & 0.9--0.99 & 0.9 default; 0.99 for noisy gradients \\
\textbf{Adam} $\beta_1$ & 0.9 & Rarely needs tuning \\
\textbf{Adam} $\beta_2$ & 0.99--0.999 & Lower (0.99) for noisy/sparse; 0.999 default \\
\textbf{Adam} $\epsilon$ & $10^{-8}$ to $10^{-6}$ & Increase to $10^{-6}$ if training unstable \\
\textbf{Weight Decay} ($\lambda$) & 0.0001--0.1 & 0.01 for AdamW; 0.0005 for SGD \\
\textbf{Gradient Clipping} & 0.5--5.0 & 1.0 common for RNNs/Transformers \\
\bottomrule
\end{tabular}
\caption{Common hyperparameter ranges.}
\label{tab:hyperparam-ranges}
\end{table}

%% ============================================================
\section{Learning Rate Schedule Recipes}\label{sec:cheat-schedules}

\begin{tcolorbox}[colback=blue!3, colframe=blue!40, title={\textbf{Step Decay}}, fonttitle=\bfseries]
\textbf{Formula:} $\eta_t = \eta_0 \cdot \gamma^{\lfloor t / s \rfloor}$ where $\gamma$ is decay factor, $s$ is step size

\textbf{When:} CNNs on image classification, when you know training will plateau

\textbf{How:}
\begin{itemize}[leftmargin=*, itemsep=1pt]
    \item Start with $\eta_0 = 0.1$ for SGD
    \item Decay by $\gamma = 0.1$ (10x reduction)
    \item Typical schedule: decay at epochs 30, 60, 90 (for 100 epochs)
    \item Or decay when validation loss plateaus for 5--10 epochs
\end{itemize}

\textbf{PyTorch:} \texttt{torch.optim.lr\_scheduler.StepLR(optimizer, step\_size=30, gamma=0.1)}
\end{tcolorbox}

\vspace{0.3cm}

\begin{tcolorbox}[colback=green!3, colframe=green!40, title={\textbf{Cosine Annealing}}, fonttitle=\bfseries]
\textbf{Formula:} $\eta_t = \eta_{\min} + \frac{1}{2}(\eta_0 - \eta_{\min})\left(1 + \cos\left(\frac{\pi t}{T}\right)\right)$

\textbf{When:} State-of-the-art CNN training, competitions, when you want smooth decay

\textbf{How:}
\begin{itemize}[leftmargin=*, itemsep=1pt]
    \item Set $\eta_0$ to your peak learning rate
    \item Set $\eta_{\min} = 0$ or $10^{-6}$
    \item $T$ = total training iterations or epochs
    \item Combine with warmup for transformers
\end{itemize}

\textbf{PyTorch:} \texttt{torch.optim.lr\_scheduler.CosineAnnealingLR(optimizer, T\_max=100)}
\end{tcolorbox}

\vspace{0.3cm}

\begin{tcolorbox}[colback=orange!3, colframe=orange!40, title={\textbf{Warmup}}, fonttitle=\bfseries]
\textbf{Formula:} $\eta_t = \eta_{\text{target}} \cdot \frac{t}{T_{\text{warmup}}}$ for $t < T_{\text{warmup}}$, then normal schedule

\textbf{When:} Large batch training, transformers, Adam/AdamW, avoiding early divergence

\textbf{How:}
\begin{itemize}[leftmargin=*, itemsep=1pt]
    \item Linear warmup for first 1--10\% of training (1000--10000 steps typical)
    \item Transformers: often 4000 steps warmup
    \item Large batches: warmup steps $\propto$ batch size
    \item Combine with cosine decay or linear decay after warmup
\end{itemize}

\textbf{PyTorch:} \texttt{torch.optim.lr\_scheduler.LinearLR(optimizer, start\_factor=0.01, total\_iters=1000)}
\end{tcolorbox}

\vspace{0.3cm}

\begin{tcolorbox}[colback=purple!3, colframe=purple!40, title={\textbf{One-Cycle Policy}}, fonttitle=\bfseries]
\textbf{Formula:} LR increases from $\eta_{\min}$ to $\eta_{\max}$ (first half), decreases to $\eta_{\min}$ (second half)

\textbf{When:} Fast training, super-convergence, when you have limited compute budget

\textbf{How:}
\begin{itemize}[leftmargin=*, itemsep=1pt]
    \item Find max LR using LR range test (loss starts increasing)
    \item Set $\eta_{\max}$ slightly below divergence point
    \item $\eta_{\min} = \eta_{\max} / 10$ or $/25$
    \item Momentum: inverse cycle (0.95 $\to$ 0.85 $\to$ 0.95)
\end{itemize}

\textbf{PyTorch:} \texttt{torch.optim.lr\_scheduler.OneCycleLR(optimizer, max\_lr=0.1, total\_steps=10000)}
\end{tcolorbox}

%% ============================================================
\section{5-Second Decision Tree}\label{sec:cheat-decision-tree}

\begin{tcolorbox}[colback=gray!5, colframe=black!70, boxrule=1pt, title={\textbf{Which Optimizer Should I Use?}}, fonttitle=\large\bfseries, coltitle=black, attach boxed title to top center={yshift=-2mm}]

\small
\renewcommand{\arraystretch}{1.4}
\begin{tabularx}{\textwidth}{@{}>{\bfseries\raggedright}p{3.8cm} >{\raggedright\arraybackslash}X@{}}
\toprule
\textbf{Scenario} & \textbf{Recommendation} \\
\midrule

Transformer / BERT / GPT &
\textbf{AdamW} + linear warmup + cosine decay \newline
{\footnotesize LR: 1e-4, warmup: 10\% of steps, weight decay: 0.01} \\[0.15cm]

CNN for competition &
\textbf{SGD + Momentum} + cosine annealing \newline
{\footnotesize LR: 0.1, momentum: 0.9, weight decay: 5e-4} \\[0.15cm]

Quick prototype &
\textbf{Adam}, LR = 3e-4 \newline
{\footnotesize The ``just works'' configuration} \\[0.15cm]

Fine-tuning pretrained &
Same optimizer as pretraining, \textbf{LR $\div$ 10} \newline
{\footnotesize Typically 1e-5 to 5e-5 for transformers} \\[0.15cm]

RNN / LSTM &
\textbf{Adam} or \textbf{RMSprop} + gradient clipping \newline
{\footnotesize LR: 1e-3, clip norm: 1.0--5.0} \\[0.15cm]

GAN training &
\textbf{Adam} with $\beta_1 = 0.5$, $\beta_2 = 0.999$ \newline
{\footnotesize LR: 2e-4, consider TTUR (different LR for G/D)} \\[0.15cm]

Sparse data / embeddings &
\textbf{AdaGrad} or \textbf{Adam} \newline
{\footnotesize Adaptive methods handle varying feature frequencies} \\[0.15cm]

Reinforcement learning &
\textbf{Adam}, LR = 3e-4 \newline
{\footnotesize Lower for policy, higher for value function} \\[0.15cm]

Very noisy gradients &
\textbf{Adam} with $\beta_2 = 0.99$ or SGD with high momentum \newline
{\footnotesize Smooths gradient estimates} \\[0.15cm]

Best generalization &
\textbf{SGD + Momentum} (properly tuned) \newline
{\footnotesize More effort but often better test accuracy} \\

\bottomrule
\end{tabularx}
\end{tcolorbox}

%% ============================================================
\section{Common Failure Modes \& Fixes}\label{sec:cheat-failures}

\begin{table}[htbp]
\centering
\renewcommand{\arraystretch}{1.4}
\begin{tabularx}{\textwidth}{@{}>{\bfseries}lX@{}}
\toprule
\textbf{Symptom} & \textbf{Fix} \\
\midrule

Loss is NaN &
$\rightarrow$ Lower LR by 10x, add gradient clipping (max\_norm=1.0), check for log(0) or division by zero \\

Loss explodes to infinity &
$\rightarrow$ LR too high---reduce by 2--10x, add gradient clipping, check data for outliers \\

Loss stuck / won't decrease &
$\rightarrow$ LR too low---increase by 2--10x, or try Adam instead of SGD \\

Loss decreases then plateaus &
$\rightarrow$ Apply LR decay (step or cosine), increase model capacity, check for data issues \\

Train loss good, val loss bad &
$\rightarrow$ Overfitting---add regularization (dropout, weight decay), early stopping, data augmentation \\

Train loss bad, val loss bad &
$\rightarrow$ Underfitting---increase model size, train longer, increase LR, reduce regularization \\

Training very slow &
$\rightarrow$ LR too low, or switch from SGD to Adam for faster initial progress \\

Adam not converging well &
$\rightarrow$ Add warmup (especially for transformers), try lower $\beta_2$ (0.99), increase $\epsilon$ \\

SGD oscillating wildly &
$\rightarrow$ LR too high, add momentum (0.9), or use LR warmup \\

Gradient clipping always active &
$\rightarrow$ LR too high, model too deep without normalization, exploding activations \\

Loss spikes during training &
$\rightarrow$ Bad data batch (check preprocessing), reduce LR, increase batch size \\

Different runs give very different results &
$\rightarrow$ Set random seeds, use larger batch size, training is sensitive---grid search LR \\

Fine-tuning destroys pretrained knowledge &
$\rightarrow$ LR too high---use $\frac{1}{10}$ to $\frac{1}{100}$ of pretraining LR, freeze early layers \\

Model trains but predicts same value &
$\rightarrow$ Check loss function, class imbalance, model architecture (missing nonlinearity) \\

\bottomrule
\end{tabularx}
\caption{Common training failures and their solutions.}
\label{tab:failure-modes}
\end{table}

%% ============================================================
\section{Quick Formulas}\label{sec:cheat-formulas}

\begin{tcolorbox}[colback=gray!5, colframe=gray!50]
\textbf{Linear LR Scaling Rule} (for batch size changes):
\[
\eta_{\text{new}} = \eta_{\text{base}} \times \frac{B_{\text{new}}}{B_{\text{base}}}
\]
Example: If LR=0.1 works for batch size 256, use LR=0.4 for batch size 1024.

\vspace{0.3cm}

\textbf{Square Root LR Scaling} (alternative, sometimes more stable):
\[
\eta_{\text{new}} = \eta_{\text{base}} \times \sqrt{\frac{B_{\text{new}}}{B_{\text{base}}}}
\]

\vspace{0.3cm}

\textbf{Gradient Accumulation} (simulate larger batch):
\[
\text{Effective batch size} = \text{micro\_batch} \times \text{accumulation\_steps}
\]
Accumulate gradients over $k$ steps, then update: equivalent to batch size $k \times B$.

\vspace{0.3cm}

\textbf{Warmup Steps Rule of Thumb}:
\[
T_{\text{warmup}} \approx 0.05 \times T_{\text{total}} \quad \text{to} \quad 0.10 \times T_{\text{total}}
\]
For transformers: minimum 1000 steps, often 4000--10000.
\end{tcolorbox}

%% ============================================================
\section{Framework Quick Reference}\label{sec:cheat-code}

\textit{Note: All code snippets below assume standard imports: \texttt{import torch}, \texttt{import numpy as np}, and relevant \texttt{torch.optim} classes. For production use, add explicit imports at the top of your scripts.}

\subsection{PyTorch}

\begin{tcolorbox}[colback=white, colframe=black!50, boxrule=0.5pt]
\begin{verbatim}
# Required imports
import torch
import torch.nn as nn
import torch.optim as optim
from torch.optim.lr_scheduler import StepLR, CosineAnnealingLR, LambdaLR

# SGD with momentum and weight decay
optimizer = torch.optim.SGD(model.parameters(),
    lr=0.1, momentum=0.9, weight_decay=5e-4)

# Adam (quick prototyping)
optimizer = torch.optim.Adam(model.parameters(), lr=3e-4)

# AdamW (transformers)
optimizer = torch.optim.AdamW(model.parameters(),
    lr=1e-4, weight_decay=0.01)

# Cosine annealing scheduler
scheduler = torch.optim.lr_scheduler.CosineAnnealingLR(
    optimizer, T_max=100)

# Warmup + cosine (manual or use transformers library)
from transformers import get_cosine_schedule_with_warmup
scheduler = get_cosine_schedule_with_warmup(
    optimizer, num_warmup_steps=1000, num_training_steps=10000)

# Gradient clipping
torch.nn.utils.clip_grad_norm_(model.parameters(), max_norm=1.0)
\end{verbatim}
\end{tcolorbox}

\subsection{TensorFlow/Keras}

\begin{tcolorbox}[colback=white, colframe=black!50, boxrule=0.5pt]
\begin{verbatim}
# Requires TensorFlow >= 2.8 for AdamW

# SGD with momentum
optimizer = tf.keras.optimizers.SGD(
    learning_rate=0.1, momentum=0.9)

# Adam (quick prototyping)
optimizer = tf.keras.optimizers.Adam(learning_rate=3e-4)

# AdamW (TensorFlow 2.8+)
try:
    optimizer = tf.keras.optimizers.AdamW(
        learning_rate=1e-4, weight_decay=0.01)
except AttributeError:
    # Fallback for older TensorFlow
    optimizer = tf.keras.optimizers.Adam(learning_rate=1e-4)

# Learning rate schedule
lr_schedule = tf.keras.optimizers.schedules.CosineDecay(
    initial_learning_rate=0.1, decay_steps=10000)
optimizer = tf.keras.optimizers.SGD(learning_rate=lr_schedule)

# Gradient clipping (in optimizer)
optimizer = tf.keras.optimizers.Adam(
    learning_rate=1e-3, clipnorm=1.0)
\end{verbatim}
\end{tcolorbox}

%% ============================================================
\section{Checklist: Before You Train}\label{sec:cheat-checklist}

\begin{tcolorbox}[colback=yellow!5, colframe=yellow!50!black, title={\textbf{Pre-Training Checklist}}, fonttitle=\bfseries]
\begin{enumerate}[leftmargin=*, itemsep=3pt]
    \item[$\square$] \textbf{Data}: Verified data loading, checked for NaN/Inf, confirmed labels correct
    \item[$\square$] \textbf{Model}: Tested forward pass, checked output shape, verified parameter count
    \item[$\square$] \textbf{Loss}: Confirmed loss function matches task (CE for classification, MSE for regression)
    \item[$\square$] \textbf{Optimizer}: Selected based on decision tree above
    \item[$\square$] \textbf{Learning Rate}: Started with recommended default for optimizer/task
    \item[$\square$] \textbf{Batch Size}: Chosen based on GPU memory, using gradient accumulation if needed
    \item[$\square$] \textbf{Scheduler}: Selected appropriate schedule (warmup for transformers, cosine for CNNs)
    \item[$\square$] \textbf{Regularization}: Added weight decay and/or dropout appropriate for model size
    \item[$\square$] \textbf{Logging}: Set up loss/metric tracking (TensorBoard, W\&B, etc.)
    \item[$\square$] \textbf{Checkpointing}: Configured model saving (best validation, periodic)
    \item[$\square$] \textbf{Sanity Check}: Verified model can overfit a single batch (loss $\to$ 0)
\end{enumerate}
\end{tcolorbox}

\vspace{0.5cm}

\begin{center}
\textit{``In doubt? Start with Adam at 3e-4. Optimize later.''}
\end{center}

\end{document}
